\section{Introduction}
\label{sec:intro}

Lightweight block ciphers protect billions of IoT sensors, RFID tags, and embedded controllers where silicon area and energy are scarce~\citep{bogdanov2007present,beaulieu2015simon}.
The three dominant design paradigms---ARX (Add-Rotate-XOR), Feistel networks, and Substitution-Permutation Networks (SPN)---each make distinct structural trade-offs between diffusion speed, parallelism, and hardware cost.
Whether one paradigm is systematically more vulnerable to a given class of attack than another is a question with direct implications for cipher selection in constrained environments.

Classical differential cryptanalysis~\citep{biham1991differential} answers this question by exhaustive search over differential trails, an approach that becomes computationally intractable as round counts grow.
Automated methods based on MILP~\citep{mouha2012milp} extend the reach of classical analysis, but they optimize within a fixed algebraic framework and cannot discover features that fall outside the differential distribution table (DDT).

In 2019, Gohr~\citep{gohr2019} demonstrated that a residual neural network, trained as a binary classifier on ciphertext pairs, could distinguish 8-round \speck from random---surpassing the best known classical distinguisher.
Since then, the field has expanded rapidly: a recent systematic survey~\citep{gerault2024sok} catalogs 66 peer-reviewed papers that train neural differential distinguishers across 24 primitives, using 15+ distinct architectures and data volumes spanning five orders of magnitude.
Bellini et al.~\citep{bellini2023dbitnet} introduced a cipher-agnostic pipeline achieving 9-round \present and 11-round \simon distinguishers, while Bao et al.~\citep{bao2022enhancing} pushed \simon key recovery to 15 rounds using neutral bits.
Yet, as Gérault et al.~\citep{gerault2024sok} explicitly note, ``the absence of systematic knowledge organization has fostered redundant research efforts,'' and no study compares transfer behavior across cipher families under controlled experimental conditions.

Our key insight is that transfer polarity---whether a distinguisher trained on one round count helps or hurts on another---reveals what the network has actually learned.
Gohr et al.~\citep{gohr2022assessment} proved that SPN and Feistel distinguishers can only learn differential features under independent round keys, while ARX distinguishers learn additional non-differential features.
If SPN features are purely differential, they should compose hierarchically across rounds (positive transfer); if ARX features include non-differential components, they may be round-specific (anti-transfer).
Our experiments confirm exactly this prediction and extend it to Feistel ciphers, where we observe anti-transfer despite the theoretical expectation of differential-only features.

We make three contributions, each falsifiable and tied to a specific experiment:
\begin{enumerate}[leftmargin=*,itemsep=2pt]
  \item \textbf{Cross-family vulnerability ordering.} Under identical architectures and training, \simon is distinguishable 2--3 rounds deeper than \speck or \present (\Cref{sec:exp-vulnerability}), with bootstrap 95\% confidence intervals for all boundary-round accuracies.
  \item \textbf{Transfer polarity is family-dependent.} \speck (ARX) and \simon (Feistel) exhibit \emph{anti-transfer} ($\leq$48.6\%, $p < 10^{-5}$), while \present (SPN) exhibits \emph{positive transfer} ($\geq$96.9\%, $p < 10^{-6}$)~(\Cref{sec:exp-transfer}). This is the first demonstration of below-chance transfer in neural cryptanalysis.
  \item \textbf{Architecture independence.} Seven architectures (MLP, CNN, Residual CNN, Gohr-MLP, LSTM, GRU, Siamese) and Gohr's ResNet all agree within 2.6\pp on the same cipher (\Cref{sec:exp-ablation}), confirming that vulnerability ordering and transfer polarity are properties of cipher structure, not model capacity.
\end{enumerate}

\Cref{sec:related} surveys related work.
\Cref{sec:problem} formalizes the threat model and notation.
\Cref{sec:method} describes the experimental methodology.
\Cref{sec:experiments} presents the main results.
\Cref{sec:discussion} interprets the findings and acknowledges limitations.
\Cref{sec:conclusion} concludes.
